\chapter{Introduction}

Urban bus network design aims to select routes, frequency, and stop patterns based on demand, the built environment, and objectives such as reliability, efficiency, and environmental impact \cite{cederBusNetworkDesign1986}. Transit agencies now have access to wide swaths of data, including \acrfull{avl}, \acrfull{afc} derived ridership patterns, and census-based demographics \cite{nchrp08-119General, bertsimasOptimizingSchoolsStart2019}. However, few existing tools combine these datasets into a computationally manageable design framework due to the complexity of balancing equity requirements, coverage targets, fleet and labor constraints, and institutional priorities. Thus, redesign efforts still largely focus on heuristic adjustments and ad-hoc methods, which limit the ability of planners and stakeholders to reason openly about tradeoffs \cite{sussmanNewApproachTransportation2005}. Given the complexity of urban bus redesigns, this leads to the question of how best to provide a scalable and reproducible method for bus network redesigns.

In this thesis, I hope to develop a general-purpose framework for representing and redesigning bus networks. Through the monotone theory of co-design, which offers a formal process for decomposing a bus redesign problem into components that maintain relationships, this framework should be able to represent the complex interplay between different datasets and stakeholders involved in transportation planning \cite{censi2016Mathematical, fong2018Seven}. Adaptable to the varying requirements of planners, I aim to incorporate insights from diverse data sources into a unified platform, including \acrlong{osm} road network information, \acrfull{gtfs} schedules, \acrfull{od} travel pairs, \acrshort{avl}-informed travel times, and U.S. \acrfull{acs} demographic information. The co-design formalism generates explicit Pareto frontiers across various functionalities and requirements of the underlying system, revealing how competing objectives relate to one another \cite{zardiniCoDesignEnableUserFriendly2023}. These fronts can allow planners and stakeholders to jointly explore and weigh alternative optimal configurations while maintaining transparency about constraints such as fleet composition, crew assignments, and prioritization of different customer segments.

A key feature of the framework will be the ability to encode existing bus route designs in the same format used to generate new alternatives. This will make it possible to evaluate current networks directly and compare them with designs produced through a \acrfull{milp} formulation \cite{patzoldLookAheadApproachesIntegrated2017}. Transit providers can therefore identify whether the current network is efficient, equitable, or if it is dominated by other feasible configurations, then determine what improvements are structurally achievable or constrained by the design space.

I will demonstrate the utility of the general framework through a case study of school bus transit in Framingham, Massachusetts, where I will assess the current routing system and generate alternative designs based on cost, coverage, and equity priorities. I will assess where different implementations lie on the generated Pareto front, illustrating how the framework can be used to explore real-world design tradeoffs in a network restructuring process transparently, while also offering a reproducible basis for planning discussions \cite{bertsimasBusRoutingOptimization2020}.

By combining a formal representational structure with integrated data pipelines and linear optimization methods, the purpose of this thesis is to highlight how the proposed approach can help planners and operators design transit systems that are both cost-effective and socially equitable, ultimately improving access across diverse urban communities. Eventually, the goal of this framework is to promote tractability, transparency, and stakeholder-aligned tradeoff analysis.

In the next section, I present an overview of relevant literature, focusing on the origins of computational bus redesigns, equity measures in this process, and its effect on school buses. I will then explain the relevant methodology and give a description of possible findings alongside a discussion of possible findings.
